\section{Task, benchmark and metrics}
\label{sec:task}

\paragraph{Task.} The input is a generated clip $V$, a target string $s$ and a hint box $b$ marking where the text is; in our benchmark $s$ is the quoted sign text of the generation prompt and $b$ comes from OCR localisation and tracking of the generated text, and the same box is given to every method (\steer uses it as its editing mask). The output is a clip $\hat V$ of the same size and length in which the text inside $b$ reads $s$ and nothing else has changed. Unlike scene text editing, the source text is a near miss produced by the generator, and in about half of the clips it is already correct, so ``do nothing'' is often the right answer.

\paragraph{Benchmark.} We generated clips of street scenes, shop fronts and screens with VACE-14B~\cite{jiang2025vace} from 120 prompts that ask for a specific sign text on one of seven carrier materials, ten seeds each (832$\times$464, 49 frames, 15\,fps). From 1{,}178 usable clips an OCR screen selected 526 as garbled; the screen is stricter than our evaluation reader, which judges \rOrigRead\% of them correct, so the benchmark contains both repair and do-nothing cases. The clips cover 40 target strings (20 English, 20 Chinese) and seven materials: neon (93), signboard (86), paper (96), glass (76), screen (65), metal (56) and wood (54); 295 are English and 231 Chinese. Forty clips inspected during development are listed in the supplement and excluded in a separate report.

\subsection{Metrics}
\label{sec:metrics}
\textbf{Readability.} A vision-language reader, Qwen2.5-VL-7B~\cite{bai2025qwen25vl}, transcribes the text inside the hint box on 5 frames; a clip is \emph{correct} if at least 4 transcriptions match the target exactly (after normalising whitespace, case and traditional/simplified Chinese). The reader is independent of \method, whose gates use a different OCR. All compared outputs are read in one job and one software environment, and a second reader from another model family re-reads the main comparison (\cref{sec:main}). \textbf{Repairs and breakages.} We count \emph{fixed} (garbled $\to$ correct) and \emph{broken} (correct $\to$ wrong) clips and their difference, the \emph{net} repairs. \textbf{Locality.} \emph{BG PSNR} is computed outside the dilated hint box over all frames; \emph{collateral change} is the grey-level change inside the text neighbourhood but outside both the old and the new glyphs (it needs the method's glyph mask, so it is reported for \method only). \textbf{Appearance.} On three frames we compare, in CIELAB, the ink colour, ink--ground contrast, halo, edge softness and stroke shading of each output with those of the original glyphs, each normalised by its median over the evaluation run; lower is closer, values are comparable only within one run, and the unseeded ink segmentation gives a test--retest noise of \rRetestSD{} per clip, so we interpret only differences well above it. \textbf{Temporal stability.} With the carrier motion estimated once on the original and applied to every output, \emph{flicker} is the motion-compensated grey-level change inside the box, \emph{jitter} the same for gradient magnitudes, and \emph{boundary jump} the change at the chunk boundaries of a chunked model. \textbf{Human judgements} come from blind pairwise studies with the original shown between two versions in random order and separate questions per quality (supplement). Full definitions are in the supplementary material, Sec.~A.
